% ID: SEL-D0038-OB-E02
% ESTADO: validada. S039, cotejada con examen y criterios.
\textbf{a)} $f'(x)=3x(x-6)$ es positiva fuera de $[0,6]$
y negativa entre $0$ y $6$. Por sus cambios de signo,
hay un máximo relativo en $\boxed{(0,8)}$ y un mínimo relativo
en $\boxed{(6,-100)}$. Además, $f''(x)=6(x-3)$ cambia de signo
en $3$, donde $f(3)=-46$. El punto de inflexión es $\boxed{(3,-46)}$.

\textbf{b)} $f(1)=0$ y $f'(1)=-15$. La tangente es
$y=-15(x-1)$, es decir, $\boxed{y=-15x+15}$.
